\thispagestyle{empty}
\begin{center}
{\large 
МИНОБРНАУКИ РОССИЙСКОЙ ФЕДЕРАЦИИ \\ФЕДЕРАЛЬНОЕ ГОСУДАРСТВЕННОЕ  
АВТОНОМНОЕ ОБРАЗОВАТЕЛЬНОЕ УЧРЕЖДЕНИЕ \\ ВЫСШЕГО ОБРАЗОВАНИЯ \\ \vspace{7pt} \textbf{«САНКТ-ПЕТЕРБУРГСКИЙ ПОЛИТЕХНИЧЕСКИЙ УНИВЕРСИТЕТ ПЕТРА ВЕЛИКОГО»}\\ \vspace{7pt} Институт компьютерных наук и кибербезопасности\\ \vspace{4pt} Высшая школа технологий искусственного ителлекта\\ \vspace{4pt} Направление: 02.03.01 «Математика и компьютерные науки»

\vspace{100pt}
\textbf{Отчет о выполнении лабораторных работ по дисциплине}\\ \vspace{5pt} \textbf{«Методы проектирования баз данных»}
\\ \vspace{5pt} \textbf{по теме «Нереляционные возможности PostgreSQL»}

}
\end{center}

\vspace{60pt}
\noindent Выполнил студент группы 5130201/40002 \vspace{-25pt} \begin{flushright} \underline{\hspace{4cm}} Гасов А.Е. \hspace{0pt}
\end{flushright} \par

\noindent Принял доцент к.т.н.\vspace{-25pt}\begin{flushright} \underline{\hspace{4cm}} Попов С.Г. \end{flushright}

\vspace{50pt}
\begin{flushright}
    «\underline{\hspace{0.6cm}}» \underline{\hspace{2.5cm}} 2026 г.
\end{flushright}

\vspace{85pt}
\begin{center}
Санкт-Петербург \\ 2026
\end{center}

\newpage